%----------------------------------------------------------------------------
\clearpage\section{Code, tools and AI}\label{app:code}
%----------------------------------------------------------------------------
\subsection{GitHub}\label{app:code:github}
Use GitHub for:
\begin{itemize}[noitemsep]
  \item code and models;
  \item notes and drawings;
  \item artifacts such as PDFs and presentations;
  \item syncing with Overleaf (\cref{app:start:access});
  \item submitting code anonymously for double-blind review, through \href{https://anonymous.4open.science/}{Anonymous GitHub};
  \item archiving with a DOI on Zenodo (\cref{app:zenodo}).
\end{itemize}
Run the repository's recurring commands (build, test, reproduce the measurements) through \href{https://github.com/casey/just}{just}, a command runner: one \texttt{justfile} lists them, and \texttt{just reproduce} is all a reviewer has to know.

\paragraph{Recommended structure.} A layout that keeps paper, code and evidence apart:
\begin{itemize}[noitemsep]
  \item \texttt{paper/}: the Overleaf-synced paper repository, included in the code repository;
  \item \texttt{src/} and \texttt{tests/}: the code;
  \item \texttt{experiments/}: measurement scripts and their raw results;
  \item \texttt{notes/}: notes, drawings and consultation summaries;
  \item \texttt{README.md}, \texttt{LICENSE}, \texttt{.zenodo.json}, \texttt{justfile} at the root.
\end{itemize}

\paragraph{README structure.} What the project is, in two sentences; how to install it; how to reproduce each result of the paper, one command per table or figure; how to cite it (the Zenodo DOI badge); the license.
\advisortodo{Drafted, as the starter-pack note on repository and README structure was empty; align it with the code repository template once that exists.}

\subsection{Behaviour-driven tests}
For behaviour-driven development (BDD), write the expected behaviour as scenarios with \href{https://cucumber.io/}{Cucumber}; the scenarios then double as evaluation cases that a reviewer can read.

\subsection{AI tools}\label{app:code:ai}
\paragraph{Disclose.} AI use must be disclosed in the appropriate form in both the TDK paper and the thesis: the faculty's declaration (\cref{app:ai}) and citations for AI-generated facts (\cref{app:writing:ai}).

\paragraph{Tools, as of October 2026.} Advice on AI tools goes stale within months; this is what we recommend now.
\begin{itemize}
  \item Use the strongest models, as they are the best for research; we prefer Claude Fable and ChatGPT Astra. Ask us if we can help you set up a subscription.
  \item Prefer an agentic harness working in your repository over a chat in the cloud: Claude Code or Codex, ideally with a good front end such as T3 Code, GitKraken Kepler or herdr. Working in a repository gives you the best features, and you \emph{need} version control when an AI edits your files.
  \item Plugins and skills: the Matt Pocock and poetato plugins for productivity; caveman, ponytail, squeez and TRON if you want to spare tokens. We highly recommend the \href{https://www.aihero.dev/skills-grill-me}{grill-me skill} (\href{https://claudemarketplaces.com/skills/mattpocock/skills/grill-me}{marketplace page}); see also \href{https://github.com/cursor/plugins/blob/main/pstack/skills/unslop/SKILL.md}{unslop} and \href{https://github.com/DietrichGebert/ponytail/blob/main/skills/ponytail/SKILL.md}{ponytail}.
  \item Do not overconstrain your model with excess instructions: bloated \texttt{AGENTS.md} files and memories degrade the performance of any model. Prefer skills with clear, broad descriptions of when to use them.
\end{itemize}

%----------------------------------------------------------------------------
\clearpage\section{Archiving the code on Zenodo}\label{app:zenodo}
%----------------------------------------------------------------------------
% TODO Delete this appendix (the \input in content/appendices.tex) before submission.
A GitHub link can move, change or disappear; a reviewer checking \textbf{correctness} (criterion D/2) needs the exact code your results came from.
\href{https://zenodo.org/}{Zenodo}, run by CERN, archives a GitHub release permanently and gives it a DOI, so cite the archived release, not the repository.
These steps follow Zenodo's documentation as of September 2026; if the interface differs, the linked pages are authoritative.

\subsection{Prepare the repository}
\begin{enumerate}
  \item A \texttt{README.md} that says how to install the code and reproduce every measurement in \cref{chap:evaluation}.
  \item A \texttt{LICENSE} file; ask your supervisor which license to use.
  \item A \texttt{.zenodo.json} file in the repository root (\cref{lst:zenodo-json}). Zenodo takes the record's title, authors, license and links from it (\href{https://help.zenodo.org/docs/github/describe-software/zenodo-json/}{format and example}). Check that it is valid JSON before releasing.
  \item If the repository has a \texttt{CITATION.cff}, keep its fields in line with \texttt{.zenodo.json}; Zenodo uses \texttt{.zenodo.json} when both exist.
\end{enumerate}

\begin{listing}[!ht]
\begin{minted}{json}
{
  "title": "TODO Artifact for: <thesis title>",
  "description": "Code and measurement scripts for the TDK paper <thesis title>.",
  "upload_type": "software",
  "version": "1.0.0",
  "publication_date": "2026-11-01",
  "license": "MIT",
  "creators": [
    { "name": "TODO Surname, Given name", "affiliation": "Budapest University of Technology and Economics", "orcid": "0000-0000-0000-0000" }
  ],
  "keywords": ["TODO"]
}
\end{minted}
\caption{A minimal \texttt{.zenodo.json}.}\label{lst:zenodo-json}
\end{listing}

\subsection{Publish the first release}
\begin{enumerate}
  \item Sign in to \href{https://zenodo.org/}{Zenodo} with the GitHub account that owns the repository. Open the profile menu, then \emph{GitHub}; click \emph{Sync now}, find the repository, and switch its toggle on (\href{https://help.zenodo.org/docs/github/enable-repository/}{enable a repository}).
  \item Commit the prepared files. On GitHub, open \emph{Releases}, then \emph{Draft a new release}: create a new tag such as \texttt{v1.0.0} on the commit your measurements used, add release notes, and publish.
  \item Back on Zenodo's GitHub page, open the repository and wait for processing. Open the new record and check its files and metadata; a failed release shows an \emph{Errors} section (\href{https://help.zenodo.org/docs/github/archive-software/github-upload/}{archive a release}).
\end{enumerate}
If the repository cannot be connected to Zenodo, upload a prepared archive instead: \emph{+}, then \emph{New upload}, fill in the metadata, \emph{Preview}, \emph{Publish} (\href{https://help.zenodo.org/docs/deposit/create-new-upload/}{upload instructions}).

\subsection{Cite it}
Each release gets its own \emph{version DOI}; the \emph{concept DOI} stands for all versions together (\href{https://support.zenodo.org/help/en-gb/1-upload-deposit/97-what-is-doi-versioning}{DOI versioning}).
\begin{itemize}
  \item In the paper, cite the \textbf{version DOI} of the release your results came from, as the \texttt{artifact} entry in \texttt{bib/references.bib} shows, and cite it in \cref{chap:implementation}.
  \item In \texttt{README.md}, put the badge of the \textbf{concept DOI}, so it always leads to the newest version. The record page's DOI badge dialog gives the Markdown.
  \item Copy both DOIs from the Zenodo record; never guess a DOI by changing its number.
\end{itemize}

\subsection{Release a new version}
Fix the code, update \texttt{version} and \texttt{publication\_date} in \texttt{.zenodo.json}, commit, and publish a new GitHub release with a new tag (for example \texttt{v1.0.1}). Zenodo adds it to the same record as a new version with a new version DOI (\href{https://help.zenodo.org/docs/deposit/manage-versions/}{manage versions}). Never move an old tag. If the paper's results now come from the new version, cite its version DOI.

